\documentclass[journal,onecolumn]{IEEEtran}

% ── Encabezado ──────────────────────────────────────────────
\usepackage{fancyhdr}
\pagestyle{fancy}
\fancyhf{}
\renewcommand{\headrulewidth}{0.4pt}
\fancyhead[L]{\small Anexo E --- Declaración de uso de IA}
\fancyfoot[C]{\thepage}

% ============================================================
% IDIOMA Y CODIFICACIÓN  (mover ANTES de csquotes)
% ============================================================
\usepackage[utf8]{inputenc}
\usepackage[T1]{fontenc}
\usepackage[spanish,es-tabla]{babel}

% ============================================================
% CITAS Y TEXTO
% ============================================================
\usepackage{csquotes}

% ============================================================
% ESPACIADO Y LISTAS
% ============================================================
\usepackage{parskip}
\setlength{\parskip}{4pt}
\usepackage{enumitem}

% ============================================================
% IDIOMA Y CODIFICACIÓN
% ============================================================
\usepackage[spanish,es-tabla]{babel}
\usepackage[utf8]{inputenc}
\usepackage[T1]{fontenc}

% ============================================================
% PAPEL Y GEOMETRÍA
% ============================================================
\usepackage{geometry}
\geometry{a4paper, top=2.5cm, bottom=2.5cm, left=2.5cm, right=2.5cm}

% ============================================================
% IMÁGENES
% ============================================================
\usepackage{graphicx}
\graphicspath{{Figuras/}}

% ============================================================
% TABLAS
% ============================================================
\usepackage{booktabs}
\usepackage{tabularx}
\usepackage{longtable}

% ============================================================
% MATEMÁTICAS
% ============================================================
\usepackage{amsmath}
\usepackage{amssymb}

% ============================================================
% CÓDIGO FUENTE
% ============================================================
\usepackage{listings}

% ============================================================
% FORMATO ACADÉMICO
% ============================================================
\usepackage{microtype}
\usepackage{float}

% ============================================================
% CAPTIONS (TÍTULOS DE FIGURAS/TABLAS)
% ============================================================
\usepackage{caption}
\captionsetup{font=small, labelfont=bf}

% ============================================================
% REFERENCIAS / HIPERVÍNCULOS
% ============================================================
\usepackage[hidelinks]{hyperref}

% ============================================================
% NUMERACIÓN Y NOMBRES EN ESPAÑOL
% ============================================================
\renewcommand{\thetable}{\arabic{table}}
\def\tablename{Tabla}

\renewcommand{\thesection}{\arabic{section}}
\renewcommand{\thesubsection}{\thesection.\arabic{subsection}}
\renewcommand{\thesubsubsection}{\thesubsection.\arabic{subsubsection}}

% ============================================================
% ESTILO DE TÍTULOS DE SECCIÓN
% ============================================================
\usepackage{titlesec}
\titleformat{\section}{\normalfont\Large\bfseries}{\thesection}{1em}{}
\titleformat{\subsection}{\normalfont\large\bfseries}{\thesubsection}{1em}{}
\titleformat{\subsubsection}{\normalfont\normalsize\bfseries}{\thesubsubsection}{1em}{}

% ============================================================
% CADA SECCIÓN PRINCIPAL (\section) INICIA EN PÁGINA NUEVA
% ============================================================
\usepackage{etoolbox}
\pretocmd{\section}{\clearpage}{}{}

% ============================================================
% BIBLIOGRAFÍA (BIBLATEX + BIBER)
% ============================================================
\usepackage[
backend=biber,
style=ieee,
sorting=none
]{biblatex}
\addbibresource{referencias_C1.bib}

% ============================================================
% DOCUMENTO
% ============================================================
\begin{document}

\subsection*{Declaración de uso de inteligencia artificial por sección}
\label{sec:declaracion-ia}

El equipo empleó Claude (Anthropic) como herramienta de asistencia de
inteligencia artificial durante el desarrollo de la PE5. La
Tabla~\ref{tab:declaracion-ia} documenta, por sección, el tipo de
asistencia recibida y el método con el que su contenido fue validado antes
de incorporarse al documento final.

\renewcommand{\arraystretch}{1.3}
\footnotesize
\begin{longtable}{
	|>{\raggedright\arraybackslash}p{3.2cm}
	|>{\raggedright\arraybackslash}p{1.8cm}
	|>{\raggedright\arraybackslash}p{4.3cm}
	|>{\raggedright\arraybackslash}p{5cm}|}
	\caption{Declaración de uso de inteligencia artificial por sección.}
	\label{tab:declaracion-ia} \\
	\hline
	\textbf{Sección} & \textbf{Herramienta} & \textbf{Tipo de asistencia} & \textbf{Método de validación} \\
	\hline
	\endfirsthead
	\multicolumn{4}{l}{\small\itshape (continuación de la Tabla~\ref{tab:declaracion-ia})} \\
	\hline
	\textbf{Sección} & \textbf{Herramienta} & \textbf{Tipo de asistencia} & \textbf{Método de validación} \\
	\hline
	\endhead
	\hline
	\multicolumn{4}{r}{\small\itshape (continúa en la página siguiente)} \\
	\endfoot
	\hline
	\endlastfoot
	Fundamento Teórico FT.1 &
	Claude (Anthropic) &
	Redacción asistida de la sección. &
	El contenido se ancla en tres cuerpos normativos aplicados
	directamente a las Secciones 5 y 8 del informe (auditoría e
	inspección), y no a un resumen genérico de las normas. \\
	\hline
	Fundamento Teórico FT.2 &
	Claude (Anthropic) &
	Redacción asistida del bloque teórico y construcción del archivo
	\texttt{referencias.bib} correspondiente. &
	Las seis fuentes citadas (ISO/IEC/IEEE 29148:2018, ISO/IEC/IEEE
	15288:2023, ISO/IEC/IEEE 12207:2017, Pohl 2025, SWEBOK v4.0, IREB
	CPRE FL v3.3.0) se verificaron de forma individual contra su
	fuente oficial antes de citarse; el archivo \texttt{.bib} y el
	\texttt{.tex} se compilaron con el ciclo completo (\texttt{pdflatex
	→ biber → pdflatex → pdflatex}) para confirmar que las citas
	resolvían correctamente antes de entregarse. \\
	\hline
	Fundamento Teórico FT.3 &
	Claude (Anthropic) &
	Redacción asistida de la sección. &
	El marco de Gotel y Finkelstein se vincula explícitamente a las
	columnas Fuente y Estado de la traza de la matriz real del punto
	6.2, en vez de citarse de forma aislada. \\
	\hline
	Fundamento Teórico FT.4 &
	Claude (Anthropic) &
	Redacción asistida de la sección. &
	La fundamentación de requisitos de IA se ancla en los tres
	componentes reales del sistema (moderación de mensajes,
	verificación de imagen, recomendación de compatibilidad) y en su
	desarrollo posterior en la Sección 7, no en generalidades sobre
	aprendizaje automático. \\
	\hline
	Fundamento Teórico FT.5 &
	Claude (Anthropic) &
	Redacción asistida de la sección. &
	El marco de PMBOK 7.\textsuperscript{a} ed. se contrasta contra el
	proceso realmente seguido por el equipo (iteraciones por práctica,
	entregables incrementales del ERS), y no se presenta como una
	descripción teórica desconectada del proyecto. \\
	\hline
	1. Introducción &
	Claude (Anthropic) &
	Redacción asistida de la sección. &
	El problema y el contexto declarados se sustentan en hallazgos
	reales de la elicitación de PE1--PE2 (falta de transparencia de
	refugios, estafas por cobros de envío, posturas heterogéneas sobre
	privacidad médica), no en una justificación genérica del dominio. \\
	\hline
	2. Metodología de Ingeniería de Requisitos &
	Claude (Anthropic) &
	Redacción asistida a partir de evidencia provista por el equipo;
	estructuración del texto y generación del código \LaTeX. &
	Cada técnica y cifra citada (entrevista, cuestionario, MoSCoW,
	Kano, EARS, inspección Fagan, defectos hallados) se contrastó
	directamente contra los informes originales en PDF de PE1--PE4 de
	los cuatro subgrupos, extraídos de sus repositorios de GitHub.
	Contenido sin respaldo verificable se excluyó explícitamente en vez
	de incluirse por inferencia. La responsable revisó y aprobó cada
	afirmación metodológica antes de integrarla. \\
	\hline
	3. ERS/SRS Final del Sistema MundiPets (3.1--3.3) &
	Claude (Anthropic) &
	Redacción asistida; extracción y verificación de datos
	directamente del archivo fuente del ERS y del historial real de
	commits del repositorio; generación de tablas y código \LaTeX. &
	La correspondencia con ISO/IEC/IEEE 29148:2018 (3.1) se verificó
	cláusula por cláusula contra la estructura real del ERS. Los
	cambios de requisitos declarados en 3.2 se limitaron a los cinco
	casos con evidencia completa de redacción antes/después. El
	commit, la fecha y la unicidad de identificadores (3.3) se
	obtuvieron ejecutando comandos reales sobre el repositorio clonado
	(\texttt{git log}, conteo automatizado), no por estimación. \\
	\hline
	4. Modelos UML del Sistema &
	Claude (Anthropic) &
	Redacción asistida de la sección. &
	Los trece casos de uso documentados (CU-01 a CU-13) se desarrollan
	en sus tres vistas de modelado (textual, secuencia, actividad) de
	forma consistente entre sí, en vez de presentarse como diagramas
	aislados sin correspondencia. \\
	\hline
	5. Validación: Inspección, Defectos y Re-inspección &
	Claude (Anthropic) &
	Redacción asistida de la sección. &
	Los 23 defectos consolidados de la inspección quedan clasificados
	y trazados individualmente (18 corregidos directamente, 5 vía
	RFC); la re-inspección de cierre confirma 0 de los 23 defectos
	originales abierto, cifra verificable en la Tabla 20 (M6) de la
	Sección 8. \\
	\hline
	6. Gestión y Trazabilidad &
	Claude (Anthropic) &
	Redacción asistida de la sección. &
	Las diez RFC consolidadas ante el Change Control Board unificado
	de la PE5 (RFC-A a RFC-J) se documentan con su estado individual de
	incorporación a la línea base vigente del ERS v4.0, con al menos
	un caso (RFC-A) verificado como ya resuelto en el ERS al momento
	de tramitarse. \\
	\hline
	7. Requisitos de los Componentes de Inteligencia Artificial &
	Claude (Anthropic) &
	Redacción asistida de la sección. &
	Las tres fichas de componentes de IA (IA-01 compatibilidad de
	cruza, IA-02 validación de imágenes, IA-03 moderación de mensajes)
	se trazan cada una a un RF de origen documental concreto (RF-06,
	RF-17), en vez de describir componentes genéricos sin anclaje al
	ERS. \\
	\hline
	8. Métricas de Calidad del ERS &
	Claude (Anthropic) &
	Redacción asistida de la sección. &
	Las seis métricas (M1--M6) se presentan con su aritmética completa
	y visible a partir de conteos reales del ERS (por ejemplo, M1a =
	27/61 × 100 = 44,26\,\%), no como porcentajes redondos sin
	respaldo. \\
	\hline
	9. Retrospectiva del Equipo &
	Claude (Anthropic) &
	Redacción asistida de la sección. &
	La retrospectiva Start--Stop--Continue se redactó referida al
	proceso real del equipo durante PE1--PE5, y no como un formato
	genérico aplicable a cualquier proyecto. \\
	\hline
	10. Conclusión 1 --- Unidad 1 &
	Claude (Anthropic) &
	Redacción asistida de la conclusión integradora. &
	La técnica de elicitación señalada como reveladora de requisitos
	nuevos (origen de RF-26 y RF-27) y el stakeholder subrepresentado
	(Refugio de animales) se verificaron contra la tabla real de
	veinticuatro evidencias (EV-01 a EV-24) del ERS, sin inferirse de
	memoria. La responsable definió el enfoque interpretativo final y
	aprobó el razonamiento antes de integrarlo al documento. \\
	\hline
	10. Conclusión 2 --- Unidad 2 &
	Claude (Anthropic) &
	Redacción asistida de la sección. &
	La conclusión sobre la especificación del ERS se ancla en cifras
	reales de la auditoría de calidad (Sección 8), como el valor
	obtenido de M1 o M3, y no en una valoración cualitativa sin
	respaldo numérico. \\
	\hline
	10. Conclusión 3 --- Unidad 3 &
	Claude (Anthropic) &
	Redacción asistida de la sección. &
	La conclusión sobre modelado y trazabilidad se sustenta en
	hallazgos concretos de la matriz final (huérfanos, cadenas rotas)
	documentados en la Sección 6, y no en una descripción genérica del
	proceso de modelado UML. \\
	\hline
	10. Conclusión 4 --- Unidad 4 &
	Claude (Anthropic) &
	Redacción asistida de la sección. &
	La conclusión sobre validación y gestión de cambios se apoya en
	las cifras verificables de la Sección 5 (23 defectos, 0 abiertos
	tras la re-inspección) y en el registro real de RFC de la Sección
	6. \\
	\hline
	10. Conclusión 5 --- Unidad 5 &
	Claude (Anthropic) &
	Redacción asistida de la sección. &
	La conclusión de cierre integra específicamente los tres
	componentes de IA reales (IA-01, IA-02, IA-03) y las métricas
	obtenidas en la Sección 8, en vez de una síntesis general no
	anclada a los artefactos del proyecto. \\
	\hline
\end{longtable}
\normalsize


\end{document}
